\chapter{Hypoglycemia}
\index{Hypoglycemia}
\label{sec:Hypoglycemia}

\section{Overview}
Hypoglycemia is a blood glucose level below 70 mg/dL, with clinically significant hypoglycemia occurring below 54 mg/dL (Level 2) and severe hypoglycemia (Level 3) being any event requiring external assistance.
\section{Causes}
This condition happens because of an imbalance between glucose-lowering therapy and glucose supply.     
\section{Risk Factors}

\begin{itemize}[noitemsep]
  \item Genetic predisposition and family history
  \item Advancing age
  \item Lifestyle factors (diet, exercise, smoking, alcohol)
  \item Environmental exposures
  \item Underlying medical conditions
  \item Medication use
\end{itemize}
\section{Signs and Symptoms}

\subsection{Common symptoms}
\begin{itemize}[noitemsep]
  \item \item You may experience adrenergic symptoms (tremor, palpitations, sweating, anxiety) and neuroglycopenic symptoms (confusion, blurred vision, seizures, loss of consciousness). Pain or discomfort in the affected area    \item Pain or discomfort in the affected area
  \end{itemize}

\subsection{Emergency warning signs}
\begin{itemize}[noitemsep]
  \item Severe or worsening symptoms of hypoglycemia require immediate medical evaluation
\end{itemize}
\section{Progression and Stages}
The condition may progress through different stages of severity over time. Early stages often involve mild or subtle symptoms that may be overlooked. Moderate stages involve more noticeable symptoms affecting daily function. Advanced stages may involve significant impairment and complications.
\section{Complications}
It is the most common acute complication of diabetes treatment (insulin and sulfonylureas).

\begin{itemize}[noitemsep]
  \item Disease progression leading to worsening symptoms
  \item Secondary infections or injuries
  \item Side effects of treatment
  \item Reduced quality of life
  \item Emotional and psychological impact
\end{itemize}
\section{Diagnosis}
\begin{itemize}[noitemsep]
  \item Doctors diagnose this based on the Whipple triad: symptoms of hypoglycemia, low measured glucose, and resolution of symptoms with glucose correction. Treatment for mild-moderate hypoglycemia is with rapid-acting oral glucose (15--20 g)
  \item Severe hypoglycemia requires IV dextrose or glucagon injection.
  \item Comprehensive medical history and physical examination
  \item Laboratory tests to assess organ function
  \item Imaging studies to visualize affected structures
  \item Specialized testing depending on the affected system
  \item Consultation with relevant specialists
\end{itemize}
\section{Prevention}
\begin{itemize}[noitemsep]
  \item Maintain a healthy lifestyle with balanced diet and exercise
  \item Avoid known risk factors
  \item Attend regular medical check-ups for early detection
  \item Follow recommended screening guidelines
\end{itemize}
\section{Treatment Options}
Medications to manage symptoms and address underlying mechanisms Lifestyle modifications and self-care measures Physical therapy and rehabilitation Surgical intervention when indicated Regular monitoring and follow-up care Patient education and support services

\begin{itemize}[noitemsep]
  \item Medications to manage symptoms and address underlying mechanisms
  \item Lifestyle modifications and self-care measures
  \item Physical therapy and rehabilitation
  \item Surgical intervention when indicated
  \item Regular monitoring and follow-up care
\end{itemize}
\section{Living with It}
\begin{itemize}[noitemsep]
  \item Follow your treatment plan as prescribed
  \item Maintain regular follow-up with your healthcare provider
  \item Make necessary lifestyle adjustments
  \item Seek support from family, friends, or support groups
  \item Stay informed about your condition
\end{itemize}
\section{Outlook}
\begin{itemize}[noitemsep]
  \item In most cases, the outlook is excellent with prompt treatment
  \item Recurrent hypoglycemia may cause hypoglycemia unawareness, requiring adjustment of insulin regimens.
\end{itemize}
